\documentclass[12pt]{article}
\usepackage[utf8]{inputenc}
\usepackage[T1]{fontenc}
\usepackage{hyperref}
\title{Tangoa Language Resources for K-12 Curriculum Development}
\author{Compiled by Claude Sonnet 5 (claude-sonnet-5) for Andr\'as Kornai}
\date{}
\begin{document}
\maketitle
\section{Introduction}
Tangoa (ISO 639-3 tgp, Glottocode tang1347), also called Movono, is an Oceanic language spoken on Tangoa Island off southern Espiritu Santo, Vanuatu. Speaker counts differ by source. Wikipedia gives about 370 (2015) and Joshua Project gives about 1,400. Tangoa is noted for its linguolabial consonants. Only reference and archival material was found.
\section{Primary Educational Resources}
No primers, textbooks or school materials were found.
\section{Linguistic and Reference Materials}
\begin{itemize}
\item Wikipedia article on Tangoa. \url{https://en.wikipedia.org/wiki/Tangoa_language}
\item Wikipedia on the Espiritu Santo languages. \url{https://en.wikipedia.org/wiki/Espiritu_Santo_languages}
\item Catalogue entry for \emph{South Sea Languages}, vol. II, which covers Tangoan. \url{https://quaibranly.fr/en/explore-collections/base/Bibliography/action/show/notice/000089237-south-sea-languages-a-series-of-studies-on-the-languages-of-the-new-hebrides-and-other-south-sea-islands-volume-ii-tangoan-santo-malo-malekula-epi-baki-and-bierian-tanna-and-futuna}
\item Rosetta Project entry for tgp. \url{https://rosettapanglossia.longnow.org/index.php?title=Tgp}
\item Overview of the languages of Santo. \url{https://nthieberger.net/vanlangs/SANTO.HTM}
\end{itemize}
\section{Religious and Community Resources}
\begin{itemize}
\item Joshua Project profile of the Tangoa people. \url{https://joshuaproject.net/people_groups/15251}
\end{itemize}
\section{Digital and Audio Resources}
\begin{itemize}
\item PARADISEC field notes, Series VES (Vanuatu, Espiritu Santo). It holds archival Santo-language recordings and materials. \url{https://paradisec.org.au/fieldnotes/VES.htm}
\end{itemize}
\section{Curriculum Development Context}
No school programme or classroom material in Tangoa was found. The archival and academic sources above could support orthography and materials development.
\end{document}